\section{Step 41 Locate DC without centering INTRA samples}\label{h263-transform-s01}
Begin with the flat forward observation. Predict the DC value from the basis
normalization before writing a loop: each dimension contributes its constant
basis factor, giving DC eight times the constant sample. A block filled with
128 therefore has DC 1024. Its AC terms vanish because their signed basis values
cancel over a constant row or column.

Keep this distinct from an INTER residual that happens to be zero. The transform
entry accepts signed sample values, and its caller decides whether they represent
INTRA samples or prediction residuals. Do not insert an unrequested level shift
inside the transform to make the two cases look alike.

\paragraph{Observe} Flat forward output begins with 1024. Zero input has all
sixty-four outputs exactly zero in either direction.
\Needspace{8\baselineskip}
\paragraph{Check} What should happen to a constant INTRA block of value 128 here?
\begin{itemize}
  \item[A.] Subtract 128 and force all coefficients to zero.
  \item[B.] Produce DC 1024 with zero AC coefficients.
  \item[C.] Treat each input byte as an already coded coefficient.
\end{itemize}

\section{Step 42 Admit metadata, then signed samples}\label{h263-transform-s02}
Check the record pointers and typed request before reading sample memory.
The input pointer needs halfword alignment; the input stride must be even and
at least sixteen bytes. The output pointer needs sixteen-byte alignment.
Validate the full-width direction and mode fields and both reserved fields.

Then read the sixty-four signed halfwords and validate their direction-specific
ranges. You can copy admitted values into a tight private block during this scan.
A late out-of-range value must not leave a partial result in caller memory.
Return sample-domain status two separately from metadata status one.

\paragraph{Observe} Stride seventeen is rejected before samples. A forward
sample of 256 is rejected after permitted sample reads, retaining the output.
\Needspace{8\baselineskip}
\paragraph{Check} Which load preserves a negative halfword residual's meaning?
\begin{itemize}
  \item[A.] A signed halfword load.
  \item[B.] An unsigned byte load of only its first byte.
  \item[C.] A pointer load used as a coefficient.
\end{itemize}

\section{Step 43 Choose the basis for the direction}\label{h263-transform-s03}
Select the forward basis or its transpose once from the validated direction.
Keep the halfword and word table pointers paired: they contain the same numerical
constants at different element widths. The inverse uses transposed basis rows;
it does not reverse a frequency array or negate all AC coefficients.

Inspect an asymmetric row and an asymmetric coefficient impulse. Either can
expose a table-orientation error hidden by a constant block. Keep vertical and
horizontal coefficient indices explicit in your address calculations.

\paragraph{Observe} The constant forward/inverse pair checks normalization,
while the off-axis impulse and mixed pattern check orientation.
\Needspace{8\baselineskip}
\paragraph{Check} What changes in the taught inverse direction?
\begin{itemize}
  \item[A.] Every coefficient becomes unsigned.
  \item[B.] Both passes are skipped.
  \item[C.] The selected basis is transposed.
\end{itemize}

\section{Step 44 Accumulate the first signed dot products}\label{h263-transform-s04}
For each input row and each basis row, multiply eight signed halfwords and add
the products. In scalar mode use full-width products and sums. In MSA mode,
\texttt{dotp\_s.w} combines pairs into four signed words; widen and reduce those
four words before storing the row result.

The input and basis signs both matter. An unsigned dot product can produce a
large positive result from a negative residual or basis constant. The final
value fits a signed word by the width proof, but the computation still needs
all eight products and their signs.

\paragraph{Observe} A negative constant block has a negative DC coefficient.
The MSA first-pass sum agrees with the scalar sum before any descaling.
\Needspace{8\baselineskip}
\paragraph{Check} Why does the signed form of the dot product matter?
\begin{itemize}
  \item[A.] It controls the memory byte order.
  \item[B.] Both samples and basis constants can be negative.
  \item[C.] It automatically quantizes the coefficients.
\end{itemize}

\section{Step 45 Store unrounded intermediates transposed}\label{h263-transform-s05}
Store the first-pass word at transposed offset $4(8u+y)$. A row of this private
buffer then contains the eight intermediate values needed by one second-pass
dot product. Its byte stride is thirty-two, while the staged halfword input's
stride is sixteen.

Do not shift away the sixteen fractional bits yet. The second basis multiplication
will add another sixteen. A premature shift changes the two-dimensional result
and can introduce bias even when a round trip looks reasonable.

\paragraph{Observe} The intermediate buffer contains 64 signed words, or 256
bytes. Its values are not yet final coefficients or reconstructed samples.
\Needspace{8\baselineskip}
\paragraph{Check} What scale remains in each stored first-pass sum?
\begin{itemize}
  \item[A.] $2^{16}$, retained for the second pass.
  \item[B.] No scale, after every row is rounded to a byte.
  \item[C.] The source picture's display aspect ratio.
\end{itemize}

\section{Step 46 Widen the second products to doublewords}\label{h263-transform-s06}
The second pass multiplies signed intermediate words by signed word copies of
the basis constants. Use doubleword products and accumulators. In MSA mode,
two four-word vector pairs cover the eight terms; signed doubleword dot products
produce four wide partial sums in total.

Combine every partial sum before rounding. Keeping only one doubleword or using
word-result dot products loses information. Because the intermediate is already
wide, this error can be much larger than a one-unit rounding difference.

\paragraph{Observe} Each private source row is thirty-two bytes. Two bounded
sixteen-byte loads cover it exactly, with no padding or extra neighbour rows.
\Needspace{8\baselineskip}
\paragraph{Check} Which result width is required for second-pass products?
\begin{itemize}
  \item[A.] A byte, because final pixels use bytes.
  \item[B.] A halfword, because original samples used halfwords.
  \item[C.] A signed doubleword.
\end{itemize}

\section{Step 47 Apply one signed descale}\label{h263-transform-s07}
Add the positive $2^{31}$ bias to the full second-pass sum, then perform the
arithmetic shift by thirty-two bits. Keep both operations full width. Check a
negative value independently; a logical right shift has a different meaning.
The bias must not become a negative sign-extended word when loaded.

This is the taught scale-$2^{16}$ implementation's convention. If you choose
another accurate precision, update the basis, bias and shift together. The
alternative scale-$2^{15}$ implementation uses $2^{29}$ and a thirty-bit shift.
Its accuracy evidence is separate from a claim of identical bit patterns.

\paragraph{Observe} With the stated tie rule, a value of $-0.75$ rounds to $-1$
and a value of $-0.5$ rounds to zero.
\Needspace{8\baselineskip}
\paragraph{Check} What goes wrong with a logical shift of a negative accumulator?
\begin{itemize}
  \item[A.] It supplies an extra reference sample.
  \item[B.] It treats the high sign bits as a large positive value.
  \item[C.] It changes an eight-by-eight block into a sixteen-by-sixteen block.
\end{itemize}

\section{Step 48 Clip the transform result, then publish}\label{h263-transform-s08}
Clip forward coefficients to the exercise's twelve-bit range, or inverse output
to its signed nine-bit range. Store the sixty-four signed halfwords privately,
then copy the completed 128-byte output to the caller. Return successful count
64 only after publication. Keep every caller byte unchanged on rejection.

Do not clip inverse output to unsigned picture bytes here. INTER reconstruction
still needs to add prediction to the signed inverse residual before final sample
clipping. Quantization and reconstruction will make that next boundary explicit.

\paragraph{Observe} Inverse clipping can return negative values down to $-256$.
The flat inverse observation returns 128 at every position.
\Needspace{8\baselineskip}
\paragraph{Check} Why must a negative inverse residual survive this function?
\begin{itemize}
  \item[A.] It may need to reduce the later predicted sample.
  \item[B.] It is an output pointer.
  \item[C.] H.263 stores every reconstructed picture sample as a negative number.
\end{itemize}

\section{Step 49 Compare with an independent finite reference}\label{h263-transform-s09}
Run public feedback against your current source:
\begin{verbatim}
python3 drafts/h263-transform-construction/grade.py \
    my-transform.asm --cli /path/to/mars-cli.jar \
    --report transform-results.json
\end{verbatim}
The 96 cases check both directions, modes and byte orders, mixed and extreme
patterns, exact zeros, padded/wide strides and rejected metadata or sample domains.
The reference computes floating-point matrix products independently of the
provided fixed-point tables. It checks every output, source access, saved state
and guard, accepting an accurate alternative precision or private frame.
Reports require a new path.

Try both flat and worked ordinary callers. A single printed first coefficient
is useful for debugging but does not test sixty-four positions. At the return
breakpoint, inspect the whole signed halfword block and the count before the
caller replaces return registers with syscall arguments.

\paragraph{Observe} Finite checks accept both supplied accurate precisions.
They reject wrong signed products, narrow second products, incorrect scaling,
wrong intermediate placement and omitted sample-domain checks.
\Needspace{8\baselineskip}
\paragraph{Check} Why use an independent reference instead of only a round trip?
\begin{itemize}
  \item[A.] A round trip guarantees every coefficient convention is correct.
  \item[B.] The reference must copy the student's intermediate registers.
  \item[C.] A forward and inverse pair can share the same normalization error.
\end{itemize}

\section{Step 50 Pass the inverse accuracy gate}\label{h263-transform-s10}
After finite checks pass, run the complete instructor accuracy gate on the target
inverse. Retain its data hashes, rounding convention, byte order, mode and per-cohort
metrics. Compare all five limits and exact zeros. An acceptable peak error does
not establish an acceptable mean error or bias.

Measure transform work separately from the host accuracy runner. The runner
includes reference data preparation, script compilation, state edits and many
function calls. For entry-to-return instruction work, include validation, staging,
both passes, reduction, clipping and publication. Do not count only the vector
dot products and call that the cost of a whole transform.

\paragraph{Observe} All eight body/mode/byte-order combinations passed the six
prescribed cohorts. Different valid output rounding choices still need their own
accuracy check before entering the reconstruction loop.
\Needspace{8\baselineskip}
\paragraph{Check} Which claim follows from passing this gate?
\begin{itemize}
  \item[A.] The complete H.263 encoder is already interoperable.
  \item[B.] The tested inverse meets the measured Annex A accuracy requirements.
  \item[C.] Quantization is lossless for every video.
\end{itemize}
